\subsection{dyadic 相位审计指数、solenoid 高度核与半圆维公理唯一性}\label{subsec:cdim-dyadic-audit-solenoid-rigidity}
本小节在定义 \ref{def:circle-dimension}--\ref{def:half-circle-dimension} 的圆维/半圆维框架下，补齐三条互相兼容的接口：
一是 dyadic 前缀读出下的相位审计指数精确律；
二是秩一局部化群 \(\ZZ\subseteq A\subseteq\QQ\) 的高度分类与通用 solenoid 因子化；
三是半圆维在“可加 + 有限指数不变 + 归一化”公理下的唯一性。

\paragraph{dyadic 前缀读出的审计深度指数}
\begin{definition}[dyadic 前缀量化与最小审计深度]\label{def:cdim-dyadic-prefix-audit-depth}
设 \(r,d\in\NN\)，\(\TT:=\RR/\ZZ\)。
对任意 \(\Theta\in M_{d\times r}(\RR)\)，定义群同态
\[
\Phi_{\Theta}:\ZZ^r\to\TT^d,\qquad
\Phi_{\Theta}(n):=\Theta n\pmod{\ZZ^d}.
\]
对 \(L\ge 1\)，定义逐坐标 dyadic 前缀量化
\[
q_L:\TT^d\to(\ZZ/2^L\ZZ)^d,\qquad
q_L(x)_j:=\left\lfloor 2^L\widetilde x_j\right\rfloor\!\!\!\pmod{2^L},
\]
其中 \(\widetilde x_j\in[0,1)\) 为 \(x_j\in\TT\) 的代表元。
再定义立方窗口
\[
C_Q:=\{n\in\ZZ^r:\ |n|_{\infty}\le Q\}\qquad(Q\ge 1),
\]
以及最小审计深度
\[
L_{\Theta}(Q):=
\min\{L\in\NN:\ q_L\circ\Phi_{\Theta}\ \text{在}\ C_Q\ \text{上单射}\}.
\]
定义 \(d\)-相位 dyadic 审计指数
\[
\alpha_{d}^{\mathrm{pref}}(\ZZ^r):=
\inf_{\Theta\in M_{d\times r}(\RR)}
\limsup_{Q\to\infty}\frac{L_{\Theta}(Q)}{\log_2 Q}.
\]
\end{definition}

\begin{theorem}[相位审计指数的精确值]\label{thm:cdim-dyadic-prefix-audit-exact}
在定义 \ref{def:cdim-dyadic-prefix-audit-depth} 的设定下，
\[
\alpha_{d}^{\mathrm{pref}}(\ZZ^r)=\frac{r}{d}.
\]
\end{theorem}
\begin{proof}
\textbf{下界.}
对任意 \(\Theta\) 与 \(L\)，映射 \(q_L\circ\Phi_{\Theta}\) 的值域基数至多为 \(2^{Ld}\)。
若其在 \(C_Q\) 上单射，则必有
\[
2^{Ld}\ge |C_Q|=(2Q+1)^r,
\]
故
\[
L_{\Theta}(Q)\ge \frac{r}{d}\log_2(2Q+1).
\]
取 \(\limsup_{Q\to\infty}\) 得
\[
\limsup_{Q\to\infty}\frac{L_{\Theta}(Q)}{\log_2 Q}\ge \frac{r}{d},
\]
再对 \(\Theta\) 取下确界，得到
\(
\alpha_{d}^{\mathrm{pref}}(\ZZ^r)\ge r/d
\)。
\par
\textbf{上界.}
由坏可逼近系统存在性（见 \cite{Schmidt1969BadlyApproximableSystems}），存在矩阵 \(\Theta\) 与常数 \(c>0\)，使对一切 \(k\in\ZZ^r\setminus\{0\}\) 有
\[
\operatorname{dist}_{\TT^d,\infty}(\Theta k,0)\ge c\,|k|_{\infty}^{-r/d},
\]
其中
\(
\operatorname{dist}_{\TT^d,\infty}(x,y):=\min_{m\in\ZZ^d}|x-y-m|_{\infty}
\)。
对任意不同 \(n,m\in C_Q\)，令 \(k=n-m\)，则 \(|k|_{\infty}\le 2Q\)，从而
\[
\operatorname{dist}_{\TT^d,\infty}\!\bigl(\Phi_{\Theta}(n),\Phi_{\Theta}(m)\bigr)
\ge c(2Q)^{-r/d}.
\]
若取 \(L\) 满足
\[
2^{-L}<c(2Q)^{-r/d},
\]
则任意两个不同点不可能落入同一 dyadic 量化盒，故 \(q_L\circ\Phi_{\Theta}\) 在 \(C_Q\) 上单射。
于是
\[
L_{\Theta}(Q)\le \frac{r}{d}\log_2 Q+O(1),
\]
故
\[
\limsup_{Q\to\infty}\frac{L_{\Theta}(Q)}{\log_2 Q}\le \frac{r}{d}.
\]
合并上下界得证。
\end{proof}

\begin{remark}[与最小分离指数口径的一致性]\label{rem:cdim-dyadic-prefix-vs-separation}
定理 \ref{thm:cdim-dyadic-prefix-audit-exact} 与
定理 \ref{thm:xi-cdim-phase-compression-aliasing-exponent}
给出的最小分离指数在主指数上等价：
对固定 \(\Theta\)，
\[
L_{\Theta}(Q)=\log_2\!\frac{1}{\Delta_Q(\Theta)}+O(1),
\]
其中 \(\Delta_Q(\Theta)\) 为 \(C_Q\) 上像点集的最小 \(\TT^d\)-距离。
因此二者在 \(\log Q\) 归一化后给出同一极限指数 \(r/d\)。
\end{remark}

\paragraph{圆维的最小相位预算与有限纤维不变性}
\begin{definition}[带固定有限外置寄存器的对数可审计]\label{def:cdim-log-auditable-with-finite-register}
设 \(G\) 为有限生成离散交换群。
取群同态
\(
\Phi:G\to\TT^d
\)
与一个固定有限群 \(R\) 及同态
\(
\rho:G\to R
\)。
称 \((\Phi,\rho)\) 在深度函数 \(L(Q)\) 下可审计，若对某个按 \(G\) 的固定生成集定义的词长球 \(B_G(Q)\)，映射
\[
g\longmapsto \bigl(q_{L(Q)}(\Phi(g)),\rho(g)\bigr)
\]
在 \(B_G(Q)\) 上单射。
若存在常数 \(C\) 使 \(L(Q)\le C\log_2(Q+2)+O(1)\)，称 \(G\) 在 \(d\) 个相位圆上\textbf{对数可审计}。
\end{definition}

\begin{proposition}[固定有限纤维不改变相位审计指数]\label{prop:cdim-finite-fiber-does-not-change-audit-exponent}
在定义 \ref{def:cdim-log-auditable-with-finite-register} 的审计口径下，
对任意有限生成离散交换群 \(G\)、有限群 \(F\) 与 \(d\in\NN\)，
相位主指数在 \(G\) 与 \(G\oplus F\) 之间相同。
\end{proposition}
\begin{proof}
设 \(G\cong \ZZ^r\oplus T\)。
对 \(G\oplus F\) 的任何可审计方案，限制到子群 \(G\oplus\{0\}\) 仍为可审计方案，故其所需主指数不小于 \(G\) 的主指数。
反向上，给定 \(G\) 的可审计方案 \((\Phi,\rho_G)\)，把有限群 \(F\) 并入外置寄存器，取
\[
\rho_{G\oplus F}(g,f):=(\rho_G(g),f)\in R_G\times F.
\]
该扩展不改变相位量化深度，只增加与 \(Q\) 无关的常数寄存器规模，故主指数不变。
两向比较即得结论。
\end{proof}

\begin{theorem}[圆维等于最小对数可审计相位圆数]\label{thm:cdim-minimal-log-auditable-phase-circles}
设 \(G\) 为有限生成离散交换群。
则在定义 \ref{def:cdim-log-auditable-with-finite-register} 的口径下，
\[
\cdim(G)=
\min\Bigl\{d\in\NN:\ G\ \text{在}\ d\ \text{个相位圆上对数可审计}\Bigr\}.
\]
\end{theorem}
\begin{proof}
写 \(G\cong \ZZ^r\oplus T\)，其中 \(r=\cdim(G)\)（定义 \ref{def:circle-dimension} 与例 \ref{ex:circle-dimension-examples}）。
\par
若 \(d<r\)，由定理 \ref{thm:cdim-dyadic-prefix-audit-exact}，自由部分 \(\ZZ^r\) 的相位审计深度主指数为 \(r/d>1\)，故不可能用 \(O(\log Q)\) 深度实现窗口单射，进而 \(G\) 亦不可。
\par
若 \(d=r\)，定理 \ref{thm:cdim-dyadic-prefix-audit-exact} 给出 \(\ZZ^r\) 上深度 \(L(Q)\le \log_2 Q+O(1)\) 的单射审计。
再由命题 \ref{prop:cdim-finite-fiber-does-not-change-audit-exponent}，有限部分 \(T\) 仅进入常数规模外置寄存器，不改变 \(\log Q\) 主阶。
故 \(d=r\) 可实现，且最小值即 \(r=\cdim(G)\)。
\end{proof}

\paragraph{秩一 solenoid 的高度分类与通用核}
\begin{definition}[高度序列与秩一局部化群]\label{def:cdim-rank-one-height-sequence}
设 \(\mathbf h=(h_p)_{p\in\mathbb P}\)，其中 \(h_p\in\NN\cup\{0,\infty\}\)。
定义
\[
A_{\mathbf h}:=
\left\{\frac{a}{b}\in\QQ:\ v_p(b)\le h_p\ \text{对所有素数 }p\right\},
\qquad
K_{\mathbf h}:=\widehat{A_{\mathbf h}}.
\]
\end{definition}

\begin{theorem}[一维 solenoid 的高度分类]\label{thm:cdim-rank-one-solenoid-height-classification}
在定义 \ref{def:cdim-rank-one-height-sequence} 的设定下，成立：
\begin{enumerate}
  \item \(K_{\mathbf h}\) 为紧、连通、交换群，且 \(\cdim(K_{\mathbf h})=1\)；
  \item 存在自然短正合列
  \[
  0\longrightarrow P_{\mathbf h}\longrightarrow K_{\mathbf h}\xrightarrow{\ \pi\ }\TT\longrightarrow 0,
  \]
  其中
  \[
  P_{\mathbf h}\cong \prod_{p} P_p(h_p),\qquad
  P_p(h_p)\cong
  \begin{cases}
  0,& h_p=0,\\
  \ZZ/p^{h_p}\ZZ,& 1\le h_p<\infty,\\
  \ZZ_p,& h_p=\infty;
  \end{cases}
  \]
  \item 任意满足 \(\ZZ\subseteq A\subseteq\QQ\) 的子群都唯一对应某个高度序列 \(\mathbf h(A)\)，并且 \(A=A_{\mathbf h(A)}\)。
\end{enumerate}
\end{theorem}
\begin{proof}
第 (3) 条是秩一无挠子群的高度分类：对每个素数定义
\[
h_p(A):=\sup\{n\in\NN:\ p^{-n}\in A\},
\]
即可得到唯一序列。
\par
由第 (3) 条与定理 \ref{thm:cdim-rankone-phase-prime-fiber}，\(K_{\mathbf h}=\widehat{A_{\mathbf h}}\) 具有短正合列
\[
0\to \widehat{A_{\mathbf h}/\ZZ}\to \widehat{A_{\mathbf h}}\to \widehat{\ZZ}\to 0,
\qquad
\widehat{\ZZ}\cong\TT.
\]
而
\[
A_{\mathbf h}/\ZZ
\cong
\bigoplus_{p:\,h_p<\infty}\ZZ/p^{h_p}\ZZ
\ \oplus\
\bigoplus_{p:\,h_p=\infty}\ZZ(p^\infty),
\]
对偶后得到
\[
\widehat{A_{\mathbf h}/\ZZ}
\cong
\prod_{p:\,h_p<\infty}\ZZ/p^{h_p}\ZZ
\times
\prod_{p:\,h_p=\infty}\ZZ_p,
\]
即第 (2) 条。
\par
第 (1) 条中，连通性来自 \(A_{\mathbf h}\) 无挠；\(\cdim=1\) 来自其 \(\QQ\)-秩为 \(1\)（见定理 \ref{thm:cdim-rankone-phase-prime-fiber}）。
证毕。
\end{proof}

\begin{corollary}[通用 solenoid 的全素数核]\label{cor:cdim-universal-solenoid-full-prime-kernel}
若对所有素数 \(p\) 取 \(h_p=\infty\)，则 \(A_{\mathbf h}=\QQ\)，并且
\[
K_{\mathbf h}=\widehat{\QQ}.
\]
此时短正合列成为
\[
0\longrightarrow \widehat{\ZZ}
\longrightarrow \widehat{\QQ}
\longrightarrow \TT
\longrightarrow 0,
\qquad
\widehat{\ZZ}\cong \prod_{p}\ZZ_p.
\]
该表达与注 \ref{rem:pom-solenoid-extension} 的接口一致。
\end{corollary}
\begin{proof}
由定理 \ref{thm:cdim-rank-one-solenoid-height-classification} 直接代入 \(h_p=\infty\) 得证。
\end{proof}

\begin{theorem}[通用 solenoid 的唯一因子化态射]\label{thm:cdim-universal-solenoid-factorization}
考虑如下对象：
\[
(K,\pi),\qquad
K\ \text{为紧连通交换群},\ \pi:K\twoheadrightarrow \TT,
\]
并要求其对偶离散群 \(A:=\widehat K\) 为秩一无挠群，且 \(\widehat\pi:\ZZ\hookrightarrow A\)。
则存在唯一连续群同态
\[
f_K:\widehat{\QQ}\longrightarrow K
\]
满足
\[
\pi\circ f_K=\pi_{\QQ},
\]
其中 \(\pi_{\QQ}:\widehat{\QQ}\twoheadrightarrow\TT\) 为推论
\ref{cor:cdim-universal-solenoid-full-prime-kernel} 中的典范投影。
\end{theorem}
\begin{proof}
由秩一无挠条件与 \(\widehat\pi:\ZZ\hookrightarrow A\)，可将 \(A\) 识别为某个
\(\ZZ\subseteq A\subseteq\QQ\)。
包含映射 \(i_A:A\hookrightarrow\QQ\) 唯一确定。
对偶化得到唯一连续同态
\[
f_K:=\widehat{i_A}:\widehat{\QQ}\to\widehat A\cong K.
\]
对偶交换图直接给出 \(\pi\circ f_K=\pi_{\QQ}\)。
唯一性由 \(i_A\) 的唯一性与 Pontryagin 对偶反变等价给出。证毕。
\end{proof}

\paragraph{交换化控制与半圆维公理唯一性}
\begin{conjecture}[交换化秩控制相位审计主指数]\label{conj:cdim-abelianization-controls-phase-audit}
设 \(\Gamma\) 为有限生成群，记其交换化自由秩
\[
r:=\operatorname{rank}\bigl(\Gamma_{\mathrm{ab}}\bigr).
\]
在仅允许群同态 \(\Gamma\to\TT^d\) 的相位读出口径下，审计深度主指数满足
\[
\alpha_d^{\mathrm{hom}}(\Gamma)=\frac{r}{d}.
\]
即非交换结构不改变主指数 \(r/d\)，仅影响常数项与次主项；其额外可审计信息必须通过非相位外置轴承载。
\end{conjecture}

\begin{proposition}[半圆维的公理唯一性]\label{prop:cdim-plus-axiomatic-uniqueness}
设 \(\delta\) 定义在有限生成可消交换幺半群类上，并满足：
\begin{enumerate}
  \item \textbf{可加性：}\(\delta(M\oplus N)=\delta(M)+\delta(N)\)；
  \item \textbf{有限指数不变性：}若 \(M\subseteq N\) 为有限 Rees 指数嵌入，则 \(\delta(M)=\delta(N)\)；
  \item \textbf{归一化：}\(\delta(\NN)=\tfrac12\)。
\end{enumerate}
则对一切此类幺半群 \(M\) 有
\[
\delta(M)=\cdim_{+}(M).
\]
\end{proposition}
\begin{proof}
记 \(G(M)\) 为 \(M\) 的 Grothendieck 完成，并写
\[
G(M)\cong \ZZ^r\oplus F,\qquad r=\operatorname{rank}G(M).
\]
由定义 \ref{def:half-circle-dimension}，
\[
\cdim_{+}(M)=\frac12\,\cdim(G(M))=\frac{r}{2}.
\]
另一方面，有限生成可消交换幺半群在饱和化后与 \(\NN^r\) 于有限指数意义下等价；
由公理 (2) 得
\(
\delta(M)=\delta(\NN^r)
\)。
再由公理 (1)(3) 得
\[
\delta(\NN^r)=r\,\delta(\NN)=\frac{r}{2}.
\]
故 \(\delta(M)=r/2=\cdim_{+}(M)\)。证毕。
\end{proof}

\endinput

